\section{Discussion}
\label{sec:discussion}

\paragraph{Effect of the two improvements}
Note that the two improvements act on different sources of error and, therefore, have to
be assessed separately. The Lagrange basis of \cref{sec:basis} removes the
ill-conditioning \emph{of the construction}. It is exact on affine elements, and on
distorted ones it reduces the condition number by orders of magnitude relative to
monomials. The remainder stems from extrapolating the metric interpolant to the bounding
box of the element (\cref{prop:kappa}). The basis cannot change the shape functions
(\cref{prop:invariance}), and consequently it cannot change the Lebesgue constant, the
patch-test behaviour in exact arithmetic, the eigenvalues or the discrete solution beyond
round-off. The hybrid space of \cref{sec:hybrid}, in contrast, changes the \emph{space},
and with it the Lebesgue constant, the stability properties and the approximation
constant. Both improvements are needed, since the hybrid space has its own Vandermonde
matrix. In the Legendre-type basis used here, this matrix is better conditioned than that
of $\Qp_p(s,t)$ in any basis for $p\ge4$ (by up to five orders of magnitude at $p=16$),
because the complement is evaluated on the exact GLL tensor grid. Its condition number
still grows with $p$ ($\approx2\cdot10^{6}$ at $p=16$, \cref{fig:kappa}).

\paragraph{Trade-off between the two trial spaces}
For smooth solutions on meshes where $\Qp_p(s,t)$ is stable, \PM\ is
more accurate than \PH\ by a factor of $1.3$--$5.5$ in the energy norm, at the same
convergence order (\cref{sec:num-convergence}). This is due to the larger Cartesian
polynomial content of $\Qp_p(s,t)$. This advantage is, however, lost abruptly when the
Lebesgue constant of $\Qp_p(s,t)$ becomes large. On the $\delta=0.2$ meshes, \PM\ diverges
for $p\ge5$, and on all three $3\times2$ meshes it develops complex and negative
eigenvalues from $p=5$--$7$ on. Neither the degree nor the amount of distortion at which
this happens is known a priori, and the Jacobian is not an indicator either (the bent
element has $\det\matr J\equiv1$). \PH\ was stable and optimally convergent in every
non-folded case studied. In our opinion, this makes $\Tp_p$ the appropriate default for
high-order unsymmetric spectral elements, in particular for transient, modal and
frequency-domain analyses. The space $\Qp_p(s,t)$, on the other hand, remains attractive
for static analyses at low to moderate degree and distortion, where its larger Cartesian
content pays off. An element-level indicator for the latter choice is $\Lam_Q$ itself,
which can be estimated by evaluating the metric shape functions on a sampling grid in the
element. Note that the condition number of $\matr V$ is not a substitute, since it also
contains the extrapolation to the bounding box (\cref{prop:kappa}).

\paragraph{Folded elements}
Exactness for polynomial solutions is observed on folded elements for both
spaces (\cref{rem:folded}), as the patch tests on the panel meshes of
\cref{sec:num-patch} confirm. At $p=8$, \PH\ reduces the round-off level there by one
to two orders of magnitude. For non-polynomial solutions on folded meshes, however, \PH\
converges less regularly than on valid meshes, and its energy error stagnates. This is due
to the fact that the complement part of $\matr K_e$ contains $1/\det\matr J$, which is
singular on the curve $\det\matr J=0$ inside the element. A remedy would be to integrate
these terms with a rule adapted to the singular curve or to write the complement
gradients in a form that avoids $\matr J^{-1}$ (e.g.\ by multiplying the complement
functions by $\det\matr J$, which changes the space). This is out of the scope of this
contribution and left for future research. On folded meshes, \PM\ is unstable at
high order in our tests, and \PP\ converges with reduced order in the nodal error.

\paragraph{Limitations}
The following limitations of the present study have to be kept in mind:
\begin{enumerate}
\item No a priori stability estimate (discrete inf-sup condition
  \citep{Babuska1971Error}) is available for either metric space. The quasi-optimality of
  \PH\ is observed for the thick cylinder (\cref{tab:quasiopt}), but not proven, and the
  statements on stability rest on the computed spectra, condition numbers and convergence
  histories. While \cref{prop:representation} gives exact representations and upper
  bounds of the Lebesgue constants, the size of the correction operators is only
  quantified numerically.
\item The completeness degree of $\Tp_p$ is $\min(p_\xi,p_\eta)$, the same as for
  $\Qp_{p_\xi,p_\eta}(s,t)$, and neither space reproduces
  $\Pp_{\max(p_\xi,p_\eta)}$.
\item The complement contributes rational integrands to $\matr K_e$ and
  $\matr M_e$. As for curved isoparametric elements, Gauss quadrature is then not exact.
  Where $\det\matr J$ vanishes inside an element, the integrands are singular.
\item The study is restricted to linear plane elasticity on quadrilaterals with at most
  biquadratic geometry. The construction carries over to hexahedra
  ($\Pp_p(x,y,z)\oplus\Cp_p$ with the $L^2$-complement of $\Pp_p(\xi,\eta,\zeta)$ in
  $\Qp_p(\xi,\eta,\zeta)$) and to triangles (complement in a simplex basis), but it has not
  been tested there.
\end{enumerate}
